\thispagestyle{empty}
\begin{center}
{\sffamily\bfseries\LARGE\color{big} Groups}\par\vspace{4pt}
{\sffamily\large Office Hour One, rewritten for independent study}\par\vspace{10pt}
\begin{minipage}{0.86\linewidth}\small\centering
A self-contained account of the first office hour of \emph{Office Hours with a Geometric Group
Theorist} (Clay and Margalit). It follows the same route through the same ideas, in the same order,
but every sentence, picture and problem here is written from scratch, and nothing is assumed
beyond school algebra. You should be able to read it without the book beside you.
\end{minipage}
\end{center}
\vspace{10pt}

\begin{prereq}
\emph{What you need.} Functions and what it means for one to be a bijection; how to compose two
functions; sets and set notation; division with remainder. Nothing else. Matrices appear in one
short example that can be skipped, and the words \emph{ring} and \emph{field} are never needed.
\end{prereq}

\begin{bigpic}
\emph{Where this is going.} A group is usually introduced as a set carrying one operation that
obeys three rules, which makes it look like a piece of bookkeeping. The operation is traditionally
called a multiplication, but the name is only a name: it can be addition, or composing two motions,
or writing one string after another, and Section~\ref{sec:def} takes the word apart. The point of this chapter is that it
is nothing of the kind: every group is the collection of symmetries of some object, and a good
choice of object turns algebraic questions into questions about a picture. The chapter ends with
the construction that makes this precise, the \emph{Cayley graph}, and with the statement that a
group is exactly the symmetry group of its own Cayley graph.
\end{bigpic}

\vspace{6pt}
{\sffamily\bfseries\large\color{big} How to read this}\par\vspace{2pt}
The coloured boxes are not decoration and they are not optional reading.
\emph{Why?} boxes contain the argument behind a claim that was only asserted.
\emph{Watch out} boxes contain the mistake that most readers make at that exact point.
\emph{Try it} boxes are short checks to do on paper before going on; they take a minute each and
they are where the understanding actually happens.
\emph{Big picture} boxes say what a piece of machinery is for.
The pictures use one convention throughout: grey labels mark \emph{fixed positions}, and a
coloured disc at a position names \emph{the thing now sitting there}. So a disc that has drifted
away from its grey label has moved, and a picture pair before and after a move shows exactly what
the move did. Every figure is placed at the moment it is needed, and every figure is referred to
in the text.

\vspace{4pt}
\hrule
\vspace{6pt}
\tableofcontents
\vspace{6pt}
\hrule
\clearpage
